\documentclass[12pt]{article}

\usepackage{amsmath}
\usepackage[margin=1in]{geometry}
\usepackage{fancyhdr}
\usepackage{amsthm}
\newtheorem{lemma}{Lemma}
\usepackage{braket}
\usepackage{amssymb}
\usepackage{pgfplots}
\usepackage{tikz}
\usepackage{bbm}
\usepackage{bm}
\usetikzlibrary{shapes.geometric}
\usepackage{enumerate}
\usepackage[shortlabels]{enumitem}

\newcommand{\pderiv}[2][]{\frac{\partial#1}{\partial#2}}

\pagestyle{fancy}
\fancyhead[l]{Hudson Benites}
\fancyhead[c]{Day 5}
\fancyfoot[c]{\thepage}
\renewcommand{\headrulewidth}{0.2pt}
\setlength{\headheight}{15pt}
\fancyhead[r]{\today}

\newcommand{\magnitude}[1]{\left\lvert #1 \right\rvert}

\begin{document}
\begin{enumerate}[label=\Alph*., leftmargin=*]
\item \begin{enumerate}[label=.\Alph*] 
    \item To show $x$ is injective, we consider two cases:\\
          \textbf{Case 1} When $n=1$, $x$ is never injective, since we have many ways to map to a single real number.\\
          \textbf{Case 2} When $n\geq 2$, $x$ is injective if and only if the columns of $A$ are linearly independent.\\
          To show $x$ is surjective, we consider two cases:\\
          \textbf{Case 1} $n\geq 3$. In this case, $x$ is never surjective because its image has at most dimension 2.\\
          \textbf{Case 2} $n=2$. In this case, we can construct any $(a,b)$ by letting $(u,v)=A^{-1}(a,b)$, and $x$ is surjective. Note that $A$ must be invertible in this case.\\
          \textbf{Case 3} $n=1$. In this case, $x$ is surjective when $A\neq0$.
\item We compute\begin{equation*}
    \pderiv[x]{u}=\pderiv{u}A\begin{bmatrix}
        u\\v
    \end{bmatrix}=A\begin{bmatrix}
        \pderiv{u}u\\
        \pderiv{u}v
    \end{bmatrix}=A\begin{bmatrix}
        1\\0
    \end{bmatrix},
    \pderiv[x]{v}=\pderiv{v}A\begin{bmatrix}
        u\\v\end{bmatrix}=A\begin{bmatrix}
        \pderiv{v}u\\
        \pderiv{v}v
    \end{bmatrix}=A\begin{bmatrix}
        0\\1
    \end{bmatrix}.
\end{equation*}
We again have cases. When the rank of $A$ is less than 2, $\pderiv[x]{u}$ and $\pderiv[x]{v}$ are linearly dependent and $x$ is not a regular surface.
When $\text{rank } A = 2$, then $x$ is a regular surface.
\end{enumerate}
\item \begin{enumerate}[label=.\Alph*]
    \item Let $x(u_1,v_1)=x(u_2,v_2)$. Then, $({u_1}^3,v_1)=({u_2}^3,v_2)$. The function $f(x)=x^3$ is injective, and $v_1=v_2$. Therefore, $x$ is injective. \par Let $(a,b)$ be a point in $\mathbb{R}^2$. Let $u=a^{1/3}$ and $v=b$. We then have $x(u,v)=(a,b)$. Thus, $x$ is surjective.
    \item We compute\begin{equation*}
        \pderiv[x]{u}=(2u^3,0), \pderiv[v]{u}=(0,1).
    \end{equation*}
    Anytime $u=0, v=c$, $\pderiv[x]{u}=(0,0)$, and therefore is not linearly independent with $\pderiv[x]{v}$. Therefore it is not a regular surface.
\end{enumerate}
\item \begin{enumerate}[label=.\Alph*]
    \item SKETCH THIS The $x_1$ and $x_2$ components of $x$ parameterize a circle via $u$ in the $xy$-plane. The $x_3$ simply takes on the value of $v\in \mathbb{R}$. This surface extends the circle into the $z$ direction, creating a cylinder. We compute
          $\pderiv[x]{u}=(-\sin u, \cos u, 0)$ and $\pderiv[x]{v}=(0,0,1)$. By observation, these are linearly independent, and thus the cylinder is a regular surface.
    \item SKETCH THIS The shape of is similar to C.A, but the radius of the circle parameterized by $u$ is dependent on $v$. As such, as the height changes so does the radius of the circle. Thus, this forms a cone with a tip at $(0,0,0)$. We compute 
          $\pderiv[x]{u}=(-v\sin u, v\cos u, 0)$ and $\pderiv[x]{v}=(\cos u, \sin u, 1)$. Anytime $v=0$, $\pderiv[x]{u}$ is the zero vector, and therefore $x$ is not a regular surface.
    \item $x$ is not a local parameterization about $p=(1,0,1)$ because it is not injective. We can construct a local parameterization by
          restricting $U$ to $U=(0,\pi)\times(0,\infty)$. Then, we can construct an $F$ defined by $F(x_1, x_2, x_3)=(\arccos(\frac{x_1}{x_3}),x_3)$.
    \item There exists no local parameterization about $p=(0,0,0)$. Only points $x(u,0)=(0,0,0)$, and this is a line, which is no longer an open subset of $\mathbb{R}^2$.
\end{enumerate}
\end{enumerate}
\end{document}